\section{Discussion}
This study investigated whether LLM-generated personalization of POI descriptions could enhance user engagement and perceived information quality while maintaining trustworthiness and accuracy. Our findings demonstrate that AI-adapted descriptions successfully achieved these objectives across multiple evaluation dimensions, with particularly notable improvements in clarity and trustworthiness.

\subsection{Effectiveness of AI Personalization}
The results reveal that AI-generated adaptations consistently shifted participant responses from neutral positions toward positive endorsements. Across all three core metrics (significance, trustworthiness, and clarity), we observed a compression of the neutral category and expansion of the top-box ratings. This pattern suggests that personalized descriptions help users form more confident and favorable assessments of POI value. The statistical significance achieved for clarity ($p = 0.0218$) and marginal significance for trustworthiness ($p = 0.050$) provide empirical support for the effectiveness of LLM-based personalization in travel content.

The near-parity in engagement preferences (221 AI vs. 220 manual) is particularly encouraging. It indicates that AI personalization does not merely optimize for a single dimension at the expense of others but rather maintains overall engagement quality while adding customization benefits. This balanced performance is critical for practical deployment, as travel platforms must ensure that personalized content remains compelling across diverse user segments.

\subsection{Trust as a Critical Factor}
The improvement in perceived trustworthiness warrants special attention. Trust is fundamental to POI recommendations, as users must believe the information is accurate and reliable before making travel decisions. The 50\% reduction in Not at all trustworthy ratings and the substantial increase in Extremely trustworthy ratings suggest that personalization, when executed thoughtfully, can enhance rather than undermine credibility. This finding challenges potential concerns that AI-generated content might be perceived as less authentic or reliable than human-authored descriptions.

We hypothesize that this trust enhancement stems from the alignment between description content and user interests. When information resonates with a user's stated preferences and travel style, it may be perceived as more relevant and therefore more credible. Additionally, the LLM's ability to maintain factual consistency while adjusting emphasis and framing likely contributed to maintaining high trust levels.

\subsection{Clarity Through Personalization}
The statistically significant improvement in clarity represents one of the study's most compelling findings. By adapting content to match user interests and expertise levels, AI-generated descriptions appear to communicate information more effectively. This suggests that generic, one-size-fits-all descriptions may inadvertently obscure key information for certain user segments. Personalization addresses this limitation by surfacing the most relevant details for each individual.

For example, a museum description personalized for a history enthusiast might emphasize historical context and artifacts, while the same POI described for a family traveler might highlight interactive exhibits and educational programs. Both versions convey accurate information but differ in emphasis and framing, resulting in clearer communication of value propositions aligned with user motivations.

\subsection{Participant Demographics and Generalizability}
Our participant pool consisted primarily of young, educated individuals (mean age 30.94 years) with intermediate to experienced travel backgrounds. The predominance of students and engineers suggests comfort with technology and potentially higher receptiveness to AI-generated content. These demographic characteristics may have contributed to the positive reception of personalized descriptions and high comfort levels with AI-authored travel content (mean = 3.94).

However, this demographic concentration also limits generalizability. Older travelers, those less familiar with AI technologies, or individuals from different cultural backgrounds might exhibit different preferences and trust patterns. The strong representation of Indian and German participants provides some international diversity but does not capture the full spectrum of global travel audiences.

\subsection{Implications for Tourism Platforms}
These findings have practical implications for travel and tourism platforms seeking to enhance user engagement. The results suggest that LLM-based personalization can be deployed to adapt POI descriptions dynamically without sacrificing trust or perceived information quality. Platforms could collect user preferences during onboarding or infer them from browsing behavior, then use LLMs to generate tailored descriptions that highlight aspects most relevant to each user.

The high comfort level with AI-generated content (mean = 3.94) indicates that transparency about AI involvement may not be a barrier to adoption, particularly among tech-savvy user segments. However, platforms should consider how to communicate the personalization process to maintain trust and set appropriate expectations.

\subsection{Limitations}
Several limitations should be acknowledged. First, participants evaluated descriptions in a controlled study environment rather than during actual trip planning, which may not fully capture real-world decision-making contexts. Second, the study did not track whether personalized descriptions influenced actual visit intentions or behaviors beyond stated preferences. Third, the within-subjects design, while enabling direct comparisons, may have introduced demand characteristics or heightened attention to differences between versions. Fourth, the study did not examine potential negative effects of over-personalization, such as filter bubbles or reduced serendipitous discovery of POIs outside user-stated interests.
